% ========== PROCESS MANAGEMENT PRIMITIVES ==========
% Spawning external tools, flexible communication patterns, and resource control
% Source: Symphony/Content/The Minimal IDE, The Assembling

\section*{Process Management Primitives}
\addcontentsline{toc}{section}{Process Management Primitives}

\lettrine{S}{ymphony's} process management primitives provide the foundation for integrating external tools, language servers, build systems, and any other command-line utilities that extensions might need. Rather than hardcoding support for specific tools, Symphony offers flexible primitives that can adapt to any external process communication pattern.

\subsection*{Spawning External Tools}
\addcontentsline{toc}{subsection}{Spawning External Tools}

The process spawning system provides control over external process creation and management:

\begin{table}[ht]
\centering
\begin{tabular}{@{}p{0.3\textwidth}p{0.35\textwidth}p{0.3\textwidth}@{}}
\toprule
\textbf{Process Type} & \textbf{Capabilities} & \textbf{Use Cases} \\
\midrule
\textbf{Long-Running Services} & 
Persistent processes, automatic restart, health monitoring & 
Language servers, build watchers, development servers \\
\midrule
\textbf{Task Execution} & 
One-time execution, progress tracking, result capture & 
Build commands, test runners, formatters \\
\midrule
\textbf{Interactive Processes} & 
Bidirectional communication, real-time interaction & 
REPLs, debuggers, interactive shells \\
\midrule
\textbf{Background Workers} & 
Asynchronous execution, resource isolation, job queuing & 
File indexing, code analysis, batch processing \\
\bottomrule
\end{tabular}
\caption{Process Types and Management Capabilities}
\end{table}

\subsection*{Flexible Communication Patterns}
\addcontentsline{toc}{subsection}{Flexible Communication Patterns}

Symphony supports multiple communication patterns between extensions and external processes:

\begin{infobox}[title=Communication Pattern Philosophy]
Different tools require different communication patterns. A language server needs bidirectional JSON-RPC, a build tool needs simple command execution, and a debugger needs complex state management. Symphony's primitives adapt to each pattern without imposing artificial constraints.
\end{infobox}

\subsubsection*{Standard I/O Communication}

The most common pattern for tool integration:

\begin{compactlist}
    \item \textbf{Stdin Writing}: Send commands and data to processes
    \item \textbf{Stdout Reading}: Receive structured output and responses
    \item \textbf{Stderr Handling}: Capture errors and diagnostic information
    \item \textbf{Line Buffering}: Efficient line-by-line processing
\end{compactlist}

\subsubsection*{Network Communication}

For tools that prefer network-based communication:

\begin{compactlist}
    \item \textbf{TCP Sockets}: Direct network connections for high-performance tools
    \item \textbf{HTTP/REST}: Web-based APIs and services
    \item \textbf{WebSockets}: Real-time bidirectional communication
    \item \textbf{Named Pipes}: High-performance local communication on supported platforms
\end{compactlist}

\subsection*{Resource Management and Control}
\addcontentsline{toc}{subsection}{Resource Management and Control}

Symphony provides resource management for external processes:

\begin{description}[leftmargin=3cm,labelwidth=2.5cm]
    \item[\textbf{Memory Limits}] Configurable memory constraints to prevent resource exhaustion
    \item[\textbf{CPU Priority}] Process scheduling priority management
    \item[\textbf{Timeout Control}] Automatic termination of long-running or stuck processes
    \item[\textbf{Health Monitoring}] Continuous process health checking and automatic restart
\end{description}

\subsection*{Integration Benefits}
\addcontentsline{toc}{subsection}{Integration Benefits}

This approach to process management provides several key advantages:

\begin{successbox}
Extensions can integrate any external tool without requiring Symphony core modifications. The generic primitives handle the complexity of process management while providing the flexibility needed for diverse integration scenarios.
\end{successbox}

\begin{compactlist}
    \item \textbf{Tool Agnostic}: Works with any command-line tool or service
    \item \textbf{Protocol Flexible}: Supports any communication protocol or data format
    \item \textbf{Resource Safe}: Prevents external processes from affecting system stability
    \item \textbf{Development Friendly}: Simplified integration patterns for extension developers
\end{compactlist}

\begin{compactlist}
    \item \textbf{TCP Sockets}: Direct socket communication for performance
    \item \textbf{HTTP/REST}: RESTful API integration for web-based tools
    \item \textbf{WebSockets}: Real-time bidirectional communication
    \item \textbf{Named Pipes}: High-performance local communication on supported platforms
\end{compactlist}

\subsubsection*{Message Protocols}

Symphony provides protocol adapters for common message formats:

\begin{lstlisting}[language=bash, caption=Protocol Communication Example]
// JSON-RPC communication with a language server
const languageServer = await symphony.processes.spawn({
    command: 'rust-analyzer',
    protocol: 'json-rpc',
    communication: 'stdio'
});

// Send initialization request
const initResult = await languageServer.request('initialize', {
    processId: process.pid,
    rootUri: workspace.rootUri,
    capabilities: clientCapabilities
});

// Handle notifications
languageServer.onNotification('textDocument/publishDiagnostics', 
    (params) => {
        diagnostics.update(params.uri, params.diagnostics);
    }
);

// Simple command execution
const buildProcess = await symphony.processes.spawn({
    command: 'cargo',
    args: ['build'],
    capture: 'all'
});

buildProcess.onExit((code, signal) => {
    if (code === 0) {
        notifications.show('Build successful');
    } else {
        problems.show(buildProcess.stderr);
    }
});
\end{lstlisting}

\subsection*{Resource Control and Isolation}
\addcontentsline{toc}{subsection}{Resource Control and Isolation}

Symphony provides resource management to ensure external processes don't impact IDE performance:

\subsubsection*{Resource Limits}

Extensions can specify resource constraints for spawned processes:

\begin{compactlist}
    \item \textbf{Memory Limits}: Maximum memory usage per process
    \item \textbf{CPU Throttling}: CPU usage limits and priority settings
    \item \textbf{File Descriptor Limits}: Control over file and network handles
    \item \textbf{Execution Timeouts}: Automatic termination of long-running processes
\end{compactlist}

\subsubsection*{Process Isolation}

Security and stability through process isolation:

\begin{compactlist}
    \item \textbf{Sandboxing}: Restricted file system and network access
    \item \textbf{User Permissions}: Process execution under restricted user accounts
    \item \textbf{Working Directory Control}: Isolated working directories
    \item \textbf{Environment Variable Filtering}: Controlled environment variable access
\end{compactlist}

\begin{successbox}
Resource control ensures that even poorly-behaved external tools cannot compromise Symphony's performance or security. Extensions can spawn resource-intensive processes with confidence that they won't affect the overall IDE experience.
\end{successbox}

\subsection*{Lifecycle Management}
\addcontentsline{toc}{subsection}{Lifecycle Management}

Symphony provides sophisticated lifecycle management for external processes:

\subsubsection*{Automatic Process Management}

\begin{expandedlist}
    \item \textbf{Auto-Start}: Processes start automatically when extensions load
    \item \textbf{Health Monitoring}: Regular health checks with configurable intervals
    \item \textbf{Crash Recovery}: Automatic restart with exponential backoff
    \item \textbf{Graceful Shutdown}: Proper cleanup when extensions are disabled
    \item \textbf{Resource Cleanup}: Automatic cleanup of temporary files and resources
\end{expandedlist}

\subsubsection*{Process State Management}

Extensions can monitor and control process state:

\begin{lstlisting}[language=JavaScript, caption=Process Lifecycle Management]
class LanguageServerManager {
    constructor() {
        this.server = null;
        this.restartCount = 0;
        this.maxRestarts = 5;
    }
    
    async start() {
        this.server = await symphony.processes.spawn({
            command: 'rust-analyzer',
            lifecycle: {
                auto_restart: true,
                health_check: this.healthCheck.bind(this),
                on_crash: this.handleCrash.bind(this)
            }
        });
        
        await this.server.waitForReady();
        this.restartCount = 0; // Reset on successful start
    }
    
    async healthCheck() {
        try {
            await this.server.request('ping', {}, { timeout: 5000 });
            return true;
        } catch (error) {
            console.warn('Language server health check failed:', error);
            return false;
        }
    }
    
    handleCrash(exitCode, signal) {
        this.restartCount++;
        
        if (this.restartCount >= this.maxRestarts) {
            notifications.showError(
                'Language server crashed too many times. Manual restart required.'
            );
            return false; // Don't auto-restart
        }
        
        return true; // Allow auto-restart
    }
    
    async stop() {
        if (this.server) {
            await this.server.shutdown();
            this.server = null;
        }
    }
}
\end{lstlisting}

\subsubsection*{Process Coordination}

For extensions that manage multiple processes:

\begin{compactlist}
    \item \textbf{Process Groups}: Logical grouping of related processes
    \item \textbf{Dependency Management}: Start/stop processes in correct order
    \item \textbf{Inter-Process Communication}: Coordination between multiple external tools
    \item \textbf{Batch Operations}: Start, stop, or restart multiple processes together
\end{compactlist}

\begin{alertbox}
The process management primitives transform Symphony into a powerful orchestration platform that can coordinate complex toolchains while maintaining simplicity and reliability.
\end{alertbox}
